\documentclass[11pt]{article}
\usepackage[margin=0.98in]{geometry}
\usepackage{amsmath,amssymb,hyperref}
\hypersetup{colorlinks=true,urlcolor=blue}
\title{The inert-prime elliptic coefficient\\ by one finite polynomial substitution}
\author{Proof-shortening proposal for ProveIt}
\date{28 November 2026}
\begin{document}
\maketitle
\noindent\textbf{Result.}
The inert-prime unit theorem in \texttt{FrobeniusCongruence.lean}
follows from a finite polynomial identity in the residue field. This
replaces the second derivative of the translated coordinate and the
ultrametric two-step induction. The proof uses neither a formal inverse
nor a square-root construction. This is a written replacement, not a
Lean-checked patch.

\section*{1. One change of variable using an integer power}
Write \(p=2m+1\). The existing proof gives \(x=z^2+\cdots\) for the
translated elliptic \(x\)-coordinate, the equation
\[
 (x')^2=4xQ(x),\qquad Q(X)=1+AX+X^2,\qquad A=a_1^2/4,
\]
and coefficient integrality through degree \(p+1\) at an inert prime
(\(p+1<p^2\)). The coefficient \(A\) is integral since \(p\) is odd and
\(a_1\) is integral. Reduce this much of the equation to the residue
field \(k\). The reduced square equation holds modulo \(z^{p+2}\):
the derivative lost one degree, but each factor \(x'\) has order at
least one. In what follows \(x\) can be any polynomial representing that
reduced jet.

Put \(u=x/z^2\), which is a polynomial with constant term 1, and set
\(t=zu^{m+1}\). The identities
\[
 t^2=xu^p,\qquad 2tt'=x'u^p,\qquad u^p\equiv1\pmod{z^p}
\]
follow from \(2(m+1)=p+1\) and characteristic \(p\).
Square the middle identity and use the equation for \(x\). After
cancelling \(4xu^p\), which has order exactly \(2\), this gives
\[
 t'^2\equiv u^p Q(x)\equiv Q(t^2)\pmod{z^p}.
\]
Therefore the square of \(t'Q(t^2)^m\) is congruent to
\(Q(t^2)^p\equiv1\pmod{z^p}\). Its constant coefficient is 1, so its
sum with 1 is a unit modulo \(z^p\), and
\[
 t'Q(t^2)^m\equiv1\pmod{z^p}.
\]
All operations can be done with polynomials and powers of \(z\) as the
precision cutoff. Neither series composition nor its inverse is needed.

\section*{2. Extracting the coefficient}
For \(0\le j<p-1\), the polynomial \(t^jt'\) is the derivative of
\(t^{j+1}/(j+1)\). Its coefficient of \(z^{p-1}\) vanishes in
characteristic \(p\). For \(j=p-1\) that coefficient is 1 since
\(t=z+\cdots\), and for \(j>p-1\) it vanishes for degree reasons.
This is the elementary differential-coefficient invariance in this jet.

As \(x^n\equiv t^{2n}Q(t^2)^mt'\pmod{z^p}\), expand the ordinary
polynomial \(X^n Q(X)^m\) and use the preceding monomial calculation:
\[
 [z^{p-1}]x^n=[X^{m-n}]Q(X)^m\quad(0\le n\le m).
\]
Writing \(c_j=[X^j]Q^m\), the case \(n=0\) gives \(c_m=0\).
The ordinary polynomial identity \(Q(Q^m)'=mQ'Q^m\) gives
\[
 (j+1)c_{j+1}=(m-j)Ac_j+(2m-j+1)c_{j-1}.
\]
If \(m=1\), then \(c_{m-1}=c_0=1\). Otherwise, if \(c_{m-1}=0\),
descend for \(j=m-1,\ldots,1\) using \(2m-j+1=p-j\ne0\) in \(k\).
This would give \(c_0=0\), a contradiction. Thus the coefficient at
\(n=1\) is nonzero in \(k\). An integral coefficient has norm 1
exactly when its residue is nonzero, proving the original unit statement.

\section*{3. Integration and checks}
\textbf{Baseline.}
The audit pins ProveIt at \texttt{f2d6ab6c4b70ea73b2deac9e33ad966bc92ea36f}.
In the
\href{https://github.com/VladimirReshetnikov/ProveIt/blob/f2d6ab6c4b70ea73b2deac9e33ad966bc92ea36f/Analysis/GeneralizedRiemannHypothesis/Lean/grh/src/New/Derivatives/FormalGroup/FrobeniusCongruence.lean#L301}{source file},
retain the first derivative square identity, the leading term, the
bound on \(a_1\) and the integer coefficient bound through \(p+1\).
Replace the local route to \texttt{norm\_coeff\_taylorVia\_translatedX\_eq\_one};
its statement is unchanged. Audit references before deleting any helpers.

\smallskip\noindent\textbf{Lean milestones.}
Reduce the integral polynomial jet coefficient by coefficient. Prove the
square congruence at precision \(p+2\); multiplication by \(x=z^2u\)
shifts the modulus by two because \(u(0)=1\). Prove the single monomial
coefficient lemma from the power rule and apply it by linearity to
\(X^nQ(X)^m\). The only recursion is on coefficients of an ordinary
polynomial over the residue field. The quotient and truncation
interfaces, total proof length, and actual dependency savings remain
to be tested in Lean.

\smallskip\noindent\textbf{Scope.}
The source audit covers all proposal TeX and README files at the pin.
The existing finite Frobenius-jet proposal concerns the coupling normal
form and exponent signatures, not this coefficient theorem. The exact
replay independently generates jets from the second-order equation
and checks the new identities in prime and quadratic finite fields.
It covers both vanishing and nonvanishing central coefficients;
only the former satisfy the full square equation at the required
precision. The replay is not a proof by finite testing.
\end{document}
